\documentclass[10pt,a4paper]{article}
\usepackage[T1]{fontenc}
\usepackage[utf8]{inputenc}
\usepackage[czech]{babel}
\usepackage{lmodern}
\usepackage{microtype}
\usepackage{geometry}
\geometry{left=14mm,right=14mm,top=18mm,bottom=18mm}
\usepackage[table]{xcolor}
\usepackage{array,longtable,tabularx,booktabs}
\usepackage{fancyhdr,lastpage}
\usepackage{ragged2e}
\usepackage{hyperref}
\hypersetup{hidelinks,pdftitle={Rozpis výuky matematiky – 4. ročník SPŠ, 112 hodin podle tematických plánů a skutečné praxe}}
\setlength{\parindent}{0pt}
\setlength{\parskip}{5pt}
\setlength{\tabcolsep}{4pt}
\renewcommand{\arraystretch}{1.08}
\setlength{\headheight}{18pt}
\pagestyle{fancy}\fancyhf{}
\fancyhead[L]{\footnotesize Matematika -- rozpis výuky}
\fancyhead[R]{\footnotesize 4. ročník SPŠ}
\fancyfoot[L]{\footnotesize Rozpis hodin podle tematického plánu}
\fancyfoot[R]{\footnotesize \thepage/\pageref{LastPage}}
\renewcommand{\headrulewidth}{0.35pt}
\newcommand{\testr}[1]{\textbf{#1}}
\begin{document}
\begin{center}
{\LARGE\bfseries Rozpis výuky matematiky}\\[3mm]
{\large 4. ročník střední průmyslové školy}\\[3mm]
{\small Přehled plánované výuky podle tematických plánů SP\v{S} a reálných zápisů výuky}
\end{center}
\vspace{3mm}
\begin{tabularx}{\textwidth}{@{}Xrrrr@{}}\toprule
\textbf{Ročník} & \textbf{Hodin} & \textbf{Krátké testy} & \textbf{Souhrnné} & \textbf{Testů celkem}\\\midrule
4. ročník & 112 & 8 & 4 & 12\\\midrule
\textbf{Celkem} & \textbf{112} & \textbf{8} & \textbf{4} & \textbf{12}\\\bottomrule
\end{tabularx}
\vspace{4mm}
\section*{Charakter dokumentu}
\begin{itemize}\setlength{\itemsep}{2pt}
\item Dokument obsahuje rozpis výuky matematiky pro 4. ročník SP\v{S} v celkové dotaci \textbf{112 hodin}.
\item Rozpis je členěn po měsících a čtvrtletích. Každý očíslovaný řádek představuje \textbf{jednu plánovanou vyučovací hodinu}.
\item Zkratka \textbf{TP} odkazuje na odpovídající tematický okruh tematického plánu. U testů je označení uvedeno přímo v náplni hodiny.
\item V každém čtvrtletí jsou zařazeny \textbf{tři testy}: dva průběžné testy (zpravidla 15--20 minut) a jedna souhrnná písemná práce.
\end{itemize}

\section*{Z čeho rozpis vychází}
\begin{itemize}\setlength{\itemsep}{2pt}
\item Hodinová dotace vychází z tematického plánu SP\v{S} pro 4. ročník: \textbf{112 hodin}.
\item Obsah a pořadí témat zohledňují reálné zápisy výuky v dostupných třídách, zejména v 1.A (2024/25), 2.A (2025/26), 3.A (2025/26) a 3.D (2025/26).
\item Historické zápisy slouží jako podklad pro reálné pořadí témat a odhad časové náročnosti; nejde o doslovný přepis třídních knih.
\end{itemize}

\section*{Podkladové dokumenty}
\begin{itemize}\setlength{\itemsep}{2pt}
\item tematické plány SP\v{S},
\item zápisy výuky.
\end{itemize}

\section*{Odlišnosti oproti tematickým plánům}
\begin{itemize}\setlength{\itemsep}{2pt}
\item \textbf{4. ročník:} protože nejsou k dispozici úplné zápisy výuky, vychází rozpis především z tematického plánu a z již dříve připravené osnovy.
\end{itemize}
\clearpage

\section*{4. ročník – posloupnosti, pravděpodobnost, statistika a opakování}
\textit{Dotace 112 hodin. Čtvrtletní celky a tři testy v každém čtvrtletí jsou součástí hodinové dotace.}
\textit{Obsah 4. ročníku není doložen historickými zápisy; vychází z TP 2020 a dřívější osnovy.}
\vspace{2mm}
\footnotesize
\begin{longtable}{@{}>{\RaggedLeft\arraybackslash}p{9mm}>{\RaggedRight\arraybackslash}p{15mm}>{\RaggedRight\arraybackslash}p{150mm}@{}}
\toprule \textbf{Hod.} & \textbf{TP} & \textbf{Náplň vyučovací hodiny} \\\midrule\endfirsthead
\toprule \textbf{Hod.} & \textbf{TP} & \textbf{Náplň vyučovací hodiny} \\\midrule\endhead
\bottomrule\endfoot
\addlinespace[1mm]\multicolumn{3}{@{}l}{\footnotesize\textbf{1. čtvrtletí} -- 2 průběžné testy + 1 souhrnná práce}\\[0.5mm]
\addlinespace[1.1mm]\multicolumn{2}{@{}l}{\textbf{Září}} & \multicolumn{1}{r@{}}{\textbf{12 hodin}}\\
\midrule
1 & 1.1 & Posloupnost jako funkce na přirozených číslech; členy a značení. \\
2 & 1.1 & Výpis členů a vzorec pro n-tý člen; hledání předpisu. \\
3 & 1.1 & Graf posloupnosti a jeho diskrétní charakter. \\
4 & 1.1 & Rekurentní zadání posloupnosti; výpočet členů. \\
5 & 1.1 & Explicitní a rekurentní zadání – srovnání a úlohy. \\
6 & 1.1 & Vlastnosti posloupnosti: monotónnost a omezenost. \\
7 & 1.2 & Aritmetická posloupnost; diference a vztah pro n-tý člen. \\
8 & 1.2 & Slovní úlohy a procvičení aritmetické posloupnosti. \\
9 & 1.3 & Geometrická posloupnost; kvocient a vztah pro n-tý člen. \\
10 & 1.3 & Součet prvních n členů geometrické posloupnosti. \\
11 & 1.3 & Procvičení geometrické posloupnosti; rozdíl oproti aritmetické. \\
12 & T1 & \testr{Průběžný test 1 (15–20 min): Posloupnosti a jejich zadávání. Zbytek hodiny: krátký rozbor a navazující procvičování.} \\
\addlinespace[1.1mm]\multicolumn{2}{@{}l}{\textbf{Říjen}} & \multicolumn{1}{r@{}}{\textbf{16 hodin}}\\
\midrule
13 & 1.4 & Užití posloupností v praxi – pravidelný přírůstek a rozdíl. \\
14 & 1.4 & Geometrický růst a pokles v technice a přírodě. \\
15 & 1.4 & Úlohy na růst počtu výrobků, výkonu či spotřeby. \\
16 & 1.4 & Srovnávací slovní úlohy na oba typy posloupností. \\
17 & 1.5 & Základní finanční pojmy: kapitál, úroková míra, období. \\
18 & 1.5 & Jednoduché úročení; výpočet úroku a konečné částky. \\
19 & 1.5 & Složené úrokování: odvození vzorce z geometrické posloupnosti. \\
20 & 1.5 & Složené úrokování: výpočet počátečního vkladu. \\
21 & 1.5 & Složené úrokování: časové období a úroková míra. \\
22 & 1.5 & Různé délky úrokovacího období; porovnání nabídek. \\
23 & 1.5 & Pravidelné vklady a spoření – jednoduchý model. \\
24 & 1.5 & Úvěry, splácení a praktická finanční gramotnost. \\
25 & T2 & \testr{Průběžný test 2 (15–20 min): Aritmetická, geometrická posloupnost a finance. Zbytek hodiny: krátký rozbor a navazující procvičování.} \\
26 & 1.4–1.5 & Smíšené aplikační úlohy na posloupnosti a úročení. \\
27 & 1.1–1.5 & Opakování před 1. písemnou prací. \\
28 & P1 & \testr{Čtvrtletní písemná práce: Posloupnosti a finanční matematika.} \\
\addlinespace[1mm]\multicolumn{3}{@{}l}{\footnotesize\textbf{2. čtvrtletí} -- 2 průběžné testy + 1 souhrnná práce}\\[0.5mm]
\addlinespace[1.1mm]\multicolumn{2}{@{}l}{\textbf{Listopad}} & \multicolumn{1}{r@{}}{\textbf{16 hodin}}\\
\midrule
29 & 2.1 & Úvod do kombinatoriky: systematický výpis, význam pořadí a opakování. \\
30 & 2.1 & Pravidlo součtu a součinu; strom možností, více kroků výběru. \\
31 & 2.2 & Faktoriál: definice, zkracování a úpravy jednoduchých výrazů. \\
32 & 2.1 & Permutace bez opakování; odvození počtu uspořádání. \\
33 & 2.1 & Permutace s opakováním; počet různých pořadí stejných prvků. \\
34 & 2.1 & Variace s opakováním; PINy, kódy a posloupnosti znaků. \\
35 & 2.1 & Kombinace bez opakování; význam neuspořádaného výběru. \\
36 & 2.2 & Kombinační číslo a jeho výpočet pomocí faktoriálu. \\
37 & 2.2 & Vlastnosti kombinačních čísel; symetrie a Pascalův vztah. \\
38 & 2.3 & Pascalův trojúhelník a jeho souvislost s kombinačními čísly. \\
39 & 2.3 & Binomická věta: rozvoj (a+b)\textasciicircum{}n pro malé exponenty. \\
40 & T3 & \testr{Průběžný test 3 (15–20 min): Variace, permutace a kombinace. Zbytek hodiny: krátký rozbor a navazující procvičování.} \\
41 & 2.1–2.3 & Smíšené kombinatorické úlohy: rozpoznání typu a řešení. \\
42 & 2.4 & Elementární jevy a výpis výsledků; hod mincí, kostkou. \\
43 & 2.4 & Zobrazení výsledkového prostoru pomocí tabulky či stromu. \\
44 & 2.5 & Náhodný jev: jistý, nemožný, opačný; množinové vyjádření. \\
\addlinespace[1.1mm]\multicolumn{2}{@{}l}{\textbf{Prosinec}} & \multicolumn{1}{r@{}}{\textbf{12 hodin}}\\
\midrule
45 & 2.5 & Klasická pravděpodobnost, stejně pravděpodobné výsledky. \\
46 & 2.5 & Pravděpodobnost v úlohách s kostkou, mincí a kartami. \\
47 & 2.5 & Pravděpodobnost s použitím kombinatorického počítání. \\
48 & 2.6 & Sjednocení a průnik jevů; výklad pomocí diagramů. \\
49 & 2.6 & Vzorec pro pravděpodobnost sjednocení dvou jevů. \\
50 & 2.6 & Pravděpodobnost opačného jevu; aplikační úlohy. \\
51 & 2.7 & Procvičování nezávislých jevů v jednoduchých pokusech. \\
52 & 2.8 & Statistický soubor, statistická jednotka a znak. \\
53 & T4 & \testr{Průběžný test 4 (15–20 min): Kombinatorika a základní pravděpodobnost. Zbytek hodiny: krátký rozbor a navazující procvičování.} \\
54 & 2.9 & Absolutní a relativní četnost; tabulka četností. \\
55 & 2.9 & Četnostní grafy a interpretace dat – praktické příklady. \\
56 & P2 & \testr{Pololetní písemná práce: Kombinatorika a základní pravděpodobnost.} \\
\addlinespace[1mm]\multicolumn{3}{@{}l}{\footnotesize\textbf{3. čtvrtletí} -- 2 průběžné testy + 1 souhrnná práce}\\[0.5mm]
\addlinespace[1.1mm]\multicolumn{2}{@{}l}{\textbf{Leden}} & \multicolumn{1}{r@{}}{\textbf{16 hodin}}\\
\midrule
57 & 2.10 & Charakteristiky polohy; aritmetický průměr. \\
58 & 2.10 & Vážený aritmetický průměr a práce s četnostmi. \\
59 & 2.10 & Modus a medián; význam odlehlých hodnot. \\
60 & 2.10 & Porovnání průměru, mediánu a modu na datech. \\
61 & 2.10 & Rozptyl – princip a výpočet pro datový soubor. \\
62 & 2.10 & Směrodatná odchylka a její interpretace. \\
63 & 2.10 & Výpočet rozptylu a odchylky pomocí kalkulačky. \\
64 & 2.10 & Charakteristiky variability z tabulky četností. \\
65 & 2.10 & Technická statistika: měření, tolerance a rozptýlení hodnot. \\
66 & 2.10 & Zpracování jednoduchého statistického šetření. \\
67 & 2.10 & Interpretace a kritické čtení statistických souhrnů. \\
68 & 2.5–2.10 & Procvičování pravděpodobnosti a statistiky. \\
69 & 2.1–2.10 & Opakování kombinatoriky, pravděpodobnosti a statistiky. \\
70 & T5 & \testr{Průběžný test 5 (15–20 min): Pravděpodobnost jevů a jejich vztahy. Zbytek hodiny: krátký rozbor a navazující procvičování.} \\
71 & 2.1–2.10 & Rozbor 2. písemné práce a odstranění nejčastějších chyb. \\
72 & 3.1 & Opakování číselných oborů, procent a poměru. \\
\addlinespace[1.1mm]\multicolumn{2}{@{}l}{\textbf{Únor}} & \multicolumn{1}{r@{}}{\textbf{12 hodin}}\\
\midrule
73 & 3.1 & Mocniny, odmocniny a úpravy výrazů. \\
74 & 3.1 & Mnohočleny, rozklad na součin, lomené výrazy. \\
75 & 3.1 & Lineární rovnice a nerovnice. \\
76 & 3.1 & Kvadratická rovnice a diskriminant. \\
77 & 3.1 & Kvadratická funkce, grafy a vrchol paraboly. \\
78 & 3.1 & Kvadratické nerovnice a odmocninové rovnice. \\
79 & 3.1 & Funkce – definiční obor, obor hodnot, graf, vlastnosti. \\
80 & T6 & \testr{Průběžný test 6 (15–20 min): Deskriptivní statistika a charakteristiky. Zbytek hodiny: krátký rozbor a navazující procvičování.} \\
81 & 3.1 & Exponenciální a logaritmické funkce. \\
82 & 3.1 & Exponenciální a logaritmické rovnice a nerovnice. \\
83 & 3.1 & Smíšený testový blok: algebra a funkce – rozbor postupů. \\
84 & P3 & \testr{Čtvrtletní písemná práce: Pravděpodobnost a statistika.} \\
\addlinespace[1mm]\multicolumn{3}{@{}l}{\footnotesize\textbf{4. čtvrtletí} -- 2 průběžné testy + 1 souhrnná práce}\\[0.5mm]
\addlinespace[1.1mm]\multicolumn{2}{@{}l}{\textbf{Březen}} & \multicolumn{1}{r@{}}{\textbf{16 hodin}}\\
\midrule
85 & 3.1 & Orientovaný úhel; převody radiánů a stupňů. \\
86 & 3.1 & Goniometrické funkce, grafy a vlastnosti. \\
87 & 3.1 & Goniometrické vzorce a úpravy výrazů. \\
88 & 3.1 & Goniometrické rovnice a interpretace řešení. \\
89 & 3.1 & Sinová a kosinová věta – obecný trojúhelník. \\
90 & 3.1 & Planimetrie: trojúhelníky, shodnost, podobnost. \\
91 & 3.1 & Planimetrie: kružnice, úhly, obsahy a obvody. \\
92 & 3.1 & Stereometrie: povrchy a objemy těles. \\
93 & 3.1 & Analytická geometrie: body a vektory v rovině. \\
94 & 3.1 & Analytická geometrie: rovnice přímky, odchylky a vzdálenosti. \\
95 & 3.1 & Analytická geometrie: přímky a roviny v prostoru. \\
96 & 3.1 & Kuželosečky: kružnice, elipsa, hyperbola a parabola. \\
97 & T7 & \testr{Průběžný test 7 (15–20 min): Opakování funkcí a rovnic. Zbytek hodiny: krátký rozbor a navazující procvičování.} \\
98 & 3.1 & Procvičování geometrických a trigonometrických úloh. \\
99 & 3.1 & Opakování před 3. písemnou prací. \\
100 & 3.1 & Rozbor 3. písemné práce a individuální oprava typických chyb. \\
\addlinespace[1.1mm]\multicolumn{2}{@{}l}{\textbf{Duben}} & \multicolumn{1}{r@{}}{\textbf{12 hodin}}\\
\midrule
101 & 3.1 & Komplexní čísla – algebraický a goniometrický tvar. \\
102 & 3.1 & Komplexní čísla – operace, mocniny a kořeny. \\
103 & 3.1 & Posloupnosti, finanční matematika – souhrnné úlohy. \\
104 & 3.1 & Pravděpodobnost a statistika – smíšené úlohy. \\
105 & 3.1 & Maturitní úlohy: algebraické výrazy a rovnice. \\
106 & T8 & \testr{Průběžný test 8 (15–20 min): Opakování geometrie a kombinatoriky. Zbytek hodiny: krátký rozbor a navazující procvičování.} \\
107 & 3.1 & Maturitní úlohy: funkce a grafy. \\
108 & 3.1 & Maturitní úlohy: planimetrie a stereometrie. \\
109 & 3.1 & Maturitní úlohy: analytická geometrie a kuželosečky. \\
110 & 3.1 & Cvičný souhrnný didaktický test či blok úloh. \\
111 & P4 & \testr{Pololetní písemná práce: Souhrnné opakování matematiky před maturitou.} \\
112 & 3.1 & Rozbor testu, závěrečné upevnění a individuální doporučení. \\
\end{longtable}
\normalsize
\end{document}
